\documentclass{article}
\usepackage{pathfinder-paper}

\title{A Comparator Gap for Distributed CVaR Decisions}
\pathfinderpapers
  {2609.04787}{Learning-Augmented Algorithms: Guarantees, Construction Mechanisms, and System-Level Implications}
  {2609.04460}{Distributed risk-averse optimization via CVaR}

\date{28 September 2026; supervisor-clarified version 2}
\begin{document}
\maketitle

\begin{abstract}
Wang et al. optimise an average of local conditional value-at-risk
(CVaR) objectives, whereas deployment may incur the CVaR of the sum of
the same losses.  We identify the comparator term needed to transfer
their finite-time bound to that aggregate objective and give a two-agent
example in which the term is actual decision loss even at an exact local
optimum.  A fresh holdout sample supports a selective aggregate-risk
certificate, subject to a declared observation interface and a safe
fallback.  Existing work has already established that the order of risk
aggregation can change decisions; our result is a precise accounting for
this particular optimisation guarantee and deployment interface.
% ledger: #6, #8, #14, #22, #23, #27
\end{abstract}

\section{The two source results and the question}

Zhao et al.\ survey learning-augmented algorithms and explain why a
component guarantee needs an explicit downstream benchmark relation,
conditional interface guarantees, and cost accounting before it becomes
a system guarantee \cite{Q}.  Their priced quantity is a predictor
invocation.  Wang et al.\ study a different primitive: agents query
sampled local losses and communicate decisions while optimising the
average of their local CVaRs \cite{P}.  Their finite-time theorem bounds
the \emph{expected} objective gap of a weighted ergodic output; their
projection makes its membership in the decision domain exact.  Neither
fact certifies a realised output against a separate aggregate-risk
budget.  These are statements of Q and P, respectively.
% ledger: #3, #11, #13

We consider only the identity deployment interface: the decision
produced by P is used for the same agents' losses.  Let the common upper-tail mass
be \(\alpha\in(0,1)\), let \(\mathcal X\) be P's decision set, and
write
\[
\begin{aligned}
C_i(x)&=\operatorname{CVaR}_{\alpha}(J_i(x,\xi_i)),
&\mathcal C(x)&=\frac1m\sum_{i=1}^m C_i(x),\\
S(x,\xi)&=\sum_{i=1}^m J_i(x,\xi_i),
&A(x)&=\operatorname{CVaR}_{\alpha}(S(x,\xi)).
\end{aligned}
\]
Here \(\xi\) has the joint deployment law and each marginal agrees
with the corresponding local evaluation law.  This joint law is an
additional specification: P's local-oracle model gives the marginals,
but no paired aggregate-loss oracle.  Set
\(\mathcal C^*=\min_{x\in\mathcal X}\mathcal C(x)\) and
\(A^*=\min_{x\in\mathcal X}A(x)\).  For P's weighted ergodic output
\(X=\hat x_T\), assume its theorem's conditions and denote a
deterministic upper bound on its right-hand side by \(B_T\).  Thus
\[
X\in\mathcal X\quad\text{pathwise},\qquad
\mathbb E[\mathcal C(X)-\mathcal C^*]\le B_T.
\]
The expectation is over P's optimisation samples and algorithmic
randomness.  This paper adds no new finite-time estimate for \(B_T\).
% ledger: #3, #6, #13, #21, #26

\section{The aggregate comparator}

\begin{proposition}[Benchmark transfer]\label{prop:bridge}
Under the preceding common-tail and matched-loss assumptions, define
\[
D=m\mathcal C^*-A^*.
\]
Then \(D\ge0\) and
\[
\mathbb E[A(X)-A^*]\le mB_T+D.
\]
% ledger: #6, #22, #27
\end{proposition}

\begin{proof}
CVaR subadditivity, also noted in P for a common tail level, gives
\(A(x)\le m\mathcal C(x)\) for every \(x\).  Minimising both sides
gives \(A^*\le m\mathcal C^*\).  At \(X\), add and subtract
\(m\mathcal C^*\), then take expectations and apply P's bound.
% ledger: #6, #22
\end{proof}

The term \(D\) compares \emph{different optima}.  P's local
suboptimality bound does not control it; P's appendix identifies the
pointwise diversification price but does not supply this optimiser-level
comparison.  The proposition is an additive bridge of the kind Q's
composition discussion requires, without yielding Q's competitive
ratio for another downstream policy.
% ledger: #11, #22, #24

\begin{proposition}[The gap can be attained]\label{prop:example}
Let \(U>0\), \(m=2\), \(\alpha=1/2\), \(\mathcal X=[-1,1]\), and let
\(Z\) be a shared fair Bernoulli variable.  For
\[
J_1(x,Z)=UZ,\qquad
J_2(x,Z)=U\bigl((1-x)Z+x(1-Z)\bigr),
\]
the exact minimiser of \(\mathcal C\) is \(x_P=1/2\), while the
minimiser of \(A\) is \(x_A=1\).  Moreover,
\[
A(x_P)-A^*=D=U/2.
\]
% ledger: #23, #27
\end{proposition}

\begin{proof}
The local and aggregate CVaRs are
\[
\begin{aligned}
C_1(x)&=U, & C_2(x)&=U\max\{1-x,x\},\\
A(x)&=U\max\{2-x,x\}=U(2-x)
&&\text{for }x\in[-1,1].
\end{aligned}
\]
Hence \(2\mathcal C^*=3U/2\), \(A^*=U\), and the displayed
gap follows.  Both losses are bounded and affine in \(x\), as required
for P's loss class.
% ledger: #23, #27
\end{proof}

This example shows actual aggregate decision loss, beyond looseness of
a certificate at a fixed point.  It does not assert that \(D\) equals
the decision loss in general.
% ledger: #22, #23, #27

\section{A conditional deployment certificate}

The aggregate benchmark comparison is an expectation bound.  To assess
a realised \(X\), freeze it and draw \(n_i\) fresh validation losses
from each marginal law, independently of P's optimisation transcript.
Suppose \(|J_i|\le U_i\), and let \(\widehat C_i(X)\) be the empirical
CVaR of this batch.  For \(\zeta\in(0,1)\), put
\[
\epsilon_i=\frac{2U_i}{\alpha}
\sqrt{\frac{\log(2m/\zeta)}{2n_i}},
\qquad
R=\sum_{i=1}^m\bigl(\widehat C_i(X)+\epsilon_i\bigr).
\]
Conditional on the optimisation transcript, the DKW inequality and P's
CVaR stability lemma give, with validation probability at least
\(1-\zeta\), simultaneous errors
\(|\widehat C_i(X)-C_i(X)|\le\epsilon_i\).  Consequently, on this
event,
\[
A(X)\le\sum_i C_i(X)\le R
\le m\mathcal C(X)+2\sum_i\epsilon_i.
\]
Freshness matters because \(X\) was chosen using the optimisation
samples.  CVaR subadditivity requires no independence across agents in
the deployment law.
% ledger: #4, #6, #12, #13

For a fixed risk budget \(b\), accepting \(X\) only if \(R\le b\)
therefore gives
\[
\Pr\{\text{accept }X\text{ and }A(X)>b\}\le\zeta.
\]
On the validation event, the risk statement also implies
\(\Pr_{\xi}\{S(X,\xi)>R\mid X,\text{validation}\}\le\alpha\)
for a fresh deployment draw.  The confidence probability over the
audit and the tail probability over future loss are different.
% ledger: #6, #12, #14

\begin{proposition}[Selective deployment]\label{prop:selective}
Suppose a known \(x_0\in\mathcal X\) satisfies \(A(x_0)\le b\),
and suppose a known margin \(m\mathcal C^*\le b-h\) holds for
\(h>0\).  If \(2\sum_i\epsilon_i\le h/2\), deploy \(X\) when
\(R\le b\) and \(x_0\) otherwise.  Write
\(\Delta=\mathcal C(x_0)-\mathcal C^*\).  Then
\[
\begin{aligned}
\Pr\{\text{reject }X\}&\le\zeta+\frac{2mB_T}{h},\\
\Pr\{A(X_{\mathrm{deploy}})>b\}&\le\zeta,\\
\mathbb E[\mathcal C(X_{\mathrm{deploy}})-\mathcal C^*]
&\le B_T+\Delta\left(\zeta+\frac{2mB_T}{h}\right).
\end{aligned}
\]
% ledger: #8, #14
\end{proposition}

\begin{proof}
On the simultaneous validation event,
\(R\le m\mathcal C(X)+2\sum_i\epsilon_i\).  The margin and
precision conditions make \(R\le b\) whenever
\(\mathcal C(X)-\mathcal C^*\le h/(2m)\).  Markov's inequality and
P's expected gap bound prove the first claim.  The fallback is known
safe, so only a false acceptance on the failed validation event can
breach \(b\).  Finally, the deployed local gap is at most the gap at
\(X\) plus \(\Delta\) on rejection; take expectations.
% ledger: #8, #14
\end{proof}

The margin and safe fallback are additional premises, not outputs of
P.  The proposition concerns the same action and same losses, rather
than an arbitrary downstream policy.
% ledger: #8, #11, #14

\section{Observation and cost limits}

Marginal validation can be arbitrarily conservative.  If exactly one
of \(m\) agents incurs \(U\) in each equally likely outcome and
\(\alpha=1/m\), then \(\sum_i C_i=mU\) but \(A=U\).
Moreover, two agents with losses \((UZ,UZ)\) or
\((UZ,U(1-Z))\), for a fair Bernoulli \(Z\), have the same local
marginals but aggregate CVaRs \(2U\) and \(U\) at \(\alpha=1/2\).
With only separately generated marginal oracle samples, their
transcripts cannot reveal which coupling will occur in deployment.
This last statement assumes no synchronised cross-agent outcomes or
other dependence information.
% ledger: #18, #21, #26

A fresh paired audit of the actual joint loss at frozen \(X\) can
instead estimate \(A(X)\) directly.  With \(n\) independent paired
vectors, \(U_\Sigma=\sum_iU_i\), and empirical aggregate CVaR
\(\widehat A(X)\), set
\[
\epsilon_J=\frac{2U_\Sigma}{\alpha}
\sqrt{\frac{\log(2/\zeta)}{2n}},
\qquad R_J=\widehat A(X)+\epsilon_J.
\]
The same DKW argument gives
\(A(X)\le R_J\le A(X)+2\epsilon_J\) with confidence
\(1-\zeta\).  This removes the local-sum slack at the selected
decision, but does not identify \(A^*\) or remove \(D\) from
Proposition~\ref{prop:bridge}.  Paired sensing and communication are
extra interface requirements.
% ledger: #22, #26, #28

If an offline local evaluation costs \(\kappa_L\) and a paired
validation vector costs \(\kappa_J\), the declared acquisition charge
is \(\kappa_L\sum_{k<T,i}s_k^i+\kappa_Jn\).  A local audit instead
charges its \(\sum_i n_i\) evaluations.  These are accounting
conventions inspired by Q, whose own charge is for predictor calls.
They exclude forming \(X\), audit communication, and any loss from
physically playing a queried action; an application must price those
as well.  The displayed risk bounds are not cost-augmented competitive
ratios.
% ledger: #11, #12, #26, #28

For a deterministic call schedule and an all-in budget \(B\) in the
same units, the risk-only acceptance threshold must subtract the
declared acquisition charge:
\[
b=B-\kappa_L\sum_{k<T,i}s_k^i-\kappa_Jn.
\]
Any additional implementation charge must also be subtracted before
the risk filter is applied.
% ledger: #11, #28

\section{Related work and interpretation}

The distinction between risk of an aggregate and aggregation of
individual risks is already published.  Almen and Dentcheva
\cite{AlmenDentcheva} compare those evaluation orders in a distributed
system and report different decisions and less conservative aggregate
risk evaluation.  Thus the main aggregation insight and the
possibility of changed decisions are prior work.  P itself notes the
pointwise diversification price \cite{P}.  The contribution here is
the exact additive comparator term for \emph{P's} finite-time
guarantee, together with its tight two-agent instance and the
assumptions needed to audit its output in Q's interface language.
% ledger: #15, #19, #22, #23, #24, #27

CVaR confidence bounds are also established.  Liang and Luo
\cite{LiangLuo} study risk-measure confidence bounds, including CVaR,
and improve on empirical-risk-plus-radius constructions.  The DKW
holdout radius used above is a direct combination of standard
concentration with P's stability lemma, not a new estimation method.
The examples about dependence are likewise illustrations of a known
risk-aggregation obstruction, specialised to this oracle interface.
% ledger: #12, #15, #18, #19, #21

\section{Limitations and open questions}

No finite-sample filter that accepts a candidate on an all-zero safe
transcript can guarantee exact distribution-free risk feasibility for
that candidate across the stated loss class.  The ledger gives a bounded,
convex, Lipschitz single-agent example with a safe fallback: the
candidate loss is \(J(x,Z)=\max\{0,x\}Z\), with fallback \(x_0=0\)
and candidate \(x=1\).  An all-zero law makes the candidate safe;
choose \(0<b<U\). Under a law where \(Z=U\) with probability
\(p\in(\alpha b/U,\alpha)\), its CVaR is \(Up/\alpha>b\).
Every finite all-zero transcript occurs with positive probability
under both laws.  A rule accepting the candidate on that transcript
can therefore falsely accept it.  Exact membership of
\(\mathcal X\) by P's projection remains a separate fact.
% ledger: #9, #10, #13

The analysis also assumes a common tail level, matched local and
deployment losses, fresh validation, and either a safe fallback or
abstention.  A paired audit needs access to the joint deployment law.
The cost expression omits implementation-specific communication and
physical-play losses.  No bound on \(D\), no aggregate-aware training
method, and no competitive ratio for a separate downstream policy
follow from the present material.
% ledger: #11, #21, #22, #24, #26, #28

A concrete next question is whether a specified application can pay
for local training, joint observation, validation and fallback while
also bounding \(D\) tightly enough to improve on its known safe
baseline.  Q's downstream sensitivity and benchmark conditions would
still need proof if the deployed policy differs from the decision
considered here.  The current results do not establish such an
improvement.
% ledger: #24, #28

\bibliographystyle{plainurl}
\bibliography{references}
\end{document}
